%% ============================================================================
\section{Randomized Verification}
\label{app:verify}
%% ============================================================================

\begin{table}[h]
\centering\small
\setlength{\tabcolsep}{4pt}
\begin{tabular}{@{}lr@{}}
\toprule
Score (random improving edits) & lowered \\
\midrule
committed world, row-stochastic tables & \fallsPctAssertedWithpriors\% \\
Shogenji measure, Eq.~\ref{eq:mrf} & \fallsPctShogenjiSigndef\% \\
prior-expectation score (App.~\ref{app:earlier}) & \fallsPctBalanced\% \\
$\LCS$ (ours; equivalently $\Lambda$) & \fallsPctLambda\% \\
\bottomrule
\end{tabular}
\caption{Share of \thmSteps{} coherence-improving edits ($\xs=\one$) that strictly lower each
score (exact enumeration).}
\label{tab:falls}
\end{table}

The verification (Table~\ref{tab:falls}) draws \thmSeq{} networks with $2$--$8$ claims, priors uniform on
$[0.05,0.95]$ and up to $2n$ relations with random couplings and weights in $[0.02,0.98]$. It
then applies one to six random elementary improvements to each, \thmSteps{} in total. $\Lambda$
failed to rise strictly on \thmFalls{} of them. For relations beyond pairs, \karySteps{} additions
of random ternary supports and conflicts (random violated sets) to random ternary networks, all
relative to $\xs=\one$, moved $\Lambda$ in the wrong direction \karyFalls{} times. With a uniformly random committed world per
network and improvements defined relative to it, $\Lambda$ failed to rise on \xsFalls{} of
\xsSteps{} edits. For $n=3$, flat priors and weight $0.8$, every one of the \uniqPairs{} pairs $(\xs,x)$
with $x\ne\xs$ admits a single support relative to $\xs$ whose addition to the empty graph lowers
the posterior gain of $x$ (Prop.~\ref{prop:unique}; \uniqBroken{} of \uniqPairs). For the only-if direction of
Theorem~\ref{thm:necessity} with the priors held fixed, \necTrials{} random cases (a random table
whose committed cell is not its extremum, a random context, and strong singleton relations
concentrating the scope on the table's extremal cell) moved $\Lambda$ in the predicted wrong
direction in \necOk{} of \necTrials. The
released scorer agrees with an independent
enumeration to $10^{-9}$.

%% ============================================================================
\section{Posterior Marginals as Priors}
\label{app:twostage}
%% ============================================================================

A probabilistic factuality pipeline returns posterior marginals $v_i$, which fix $\xs$ by
thresholding and serve as the MRF's priors, $\pi_i=v_i$ (\S\ref{sec:setup}). The released scorer
does this by default: a value strictly inside $(0,1)$ is used as a prior, and a value of exactly $0$
or $1$ is read as a hard label and gets the flat prior $\frac12$, so label-only and probabilistic
pipelines need no configuration. Theorem~\ref{thm:mono} holds for any priors in $(0,1)$, so the
guarantee is the same as with flat priors; what changes is how much each edit counts. The upper bound of Prop.~\ref{prop:props}(v) becomes
$-\log\pi(\xs)=-\sum_i\log\pi_i^{\xs_i}(1-\pi_i)^{1-\xs_i}$, so an argument whose committed values
the priors hold only weakly has more to gain from support. A support gains $-\log(1-p_rq_r)$, most
when its violation is probable under the model, so that it rules out plausible worlds other than
$\xs$. A conflict costs $\log\frac{1-p_r}{1-p_rq_r}$, most when the rest of the model already
deems the violated configuration improbable.

Confident posteriors affect the two quantities of \S\ref{sec:score} differently. As the $v_i$
approach $0$ and $1$ at $\xs$, every violation probability $q_r$ approaches $0$ for supports and
$1$ for conflicts, both closed forms approach $0$, and $\Lambda$ vanishes for every graph:
$\pi(\xs)\to1$, so the relations have no probability left to add to $\xs$. $\LCS$ does not
tend to $\frac12$. It is the logistic of the weight of evidence
$\log P(E\mid\xs)/P(E\mid\text{not }\xs)$ of the check model (\S\ref{sec:mrf}), and both
probabilities stay bounded away from zero. The first, $\prod_r\psi_r(\xs)$, does not depend
on the priors. The second is the probability of passing all checks for a world other than $\xs$,
and as the priors sharpen it is dominated by the worlds that differ from $\xs$ in a single claim.
On Figure~\ref{fig:voyager}, priors of $0.9/0.1$, $0.99/0.01$ and $0.999/0.001$ at $\xs$ give
$\Lambda$ of $\degAOne$, $\degATwo$, $\degAThree$ for A and $\degBOne$, $\degBTwo$, $\degBThree$
for B, but $\LCS$ of $\degALCSOne$, $\degALCSTwo$, $\degALCSThree$ for A and $\degBLCSOne$,
$\degBLCSTwo$, $\degBLCSThree$ for B, against $\vyALCS$ and $\vyBLCS$ under flat priors: the scale
moves with the priors, but A stays far above B.
Theorem~\ref{thm:mono} holds throughout. Confident posteriors therefore leave $\LCS$ informative,
one of the reasons it, rather than $\Lambda$, is the reported score (\S\ref{sec:score}). Two caveats apply. Under the
scorer's default rule a posterior of exactly $1$ gets the flat prior while $0.999$ is used as is,
so the score jumps between the two (for B, $\LCS$ moves from $\degBLCSThree$ to $\vyBLCS$); a
pipeline that reports exact $0/1$ posteriors is best treated as label-only throughout. And the posteriors
should be computed without the response's claim--claim relations, as they are when the factuality
model relates claims only to retrieved evidence; otherwise those relations are counted twice.

